\newcommand{\Poly}{\mathcal{P}}
\newcommand{\TS}{\mathcal{T}}

\section{Upper Bound}
In this section, we prove that for any $\zesn$ instance, there is a solution of size only depending on $k$ such that achieves stretch at most $1$.

We use the metric space defined by the tight span defined by the distance between the terminals. Consider the $k$-dimensional vector space, we define the polytope $\Poly(T)$ to be the set of vecotrs $x$ such that for any $i$ and $j$, $x_i+x_j \ge \ell(t_i,t_j)$. The tight span $\TS(T)$ is the set of vectors $x \in \Poly(T)$ such that for any $i$, there exists a $j$ such that $x_i+x_j = \ell(t_i,t_j)$. The tight span metric space is the metric whose point set is $\TS(T)$, and the distance between two points $x$ and $y$ is defined by the $\ell_{\infty}$ norm $\norm{x-y}_{\infty}$. For any terminal $t_i$, we define the vector represent it as $x^{t_i} = (\ell(t_1,t_i),\ell(t_2,t_i),\dots,\ell(t_k,t_i))$. The following claim shows that for any vector $x$, $x_i$ is the distance between $x$ and $t_i$.

\begin{claim} 
    For any $x$ in the metric space, $\norm{x^{t_i} - x}_{\infty} = x_i$.
\end{claim}

\begin{proof}
    For any $j \neq i$, $x_i+x_j\ge \ell(t_i,t_j)$ by triangle inequality, so $x^{t_i}_j-x_j \le x_i$. Furthermore, by minimality, there exists an index $j'$ such that $x_j+x_{j'} = \ell(t_j,t_{j'})$. We have 
    $$x_j -x^{t_i}_j = \ell(t_j,t_{j'}) - x_{j'} - \ell(t_i,t_j) \le \ell(t_i,t_{j'}) - x_{j'} \le x_i$$
    by triangle inequality. Therefore, $\card{x^{t_i}_j-x_j} \le x_i$ for any $j \neq i$ and $\card{x^{t_i}_i-x_i} = x_i$, which means $\norm{x^{t_i} - x}_{\infty} = x_i$.
\end{proof}

We first project each vertex in graph to the tight span metric $\TS(T)$, and prove that the distance between each pair of vertex will only decrease; and then we define a set of steiner nodes and give a random scheme that project each point in the metric space to the steiner nodes such that any two points in the metric space will keep their distance in expectation.

\subsection{Projection of the Points}
In this section, we project each point in the graph to the metric space. Given a point $v$, suppose the distances between $v$ and each terminal are $d_1,d_2,\dots,d_k$. We construct the projection point $x=f(v)$ by the following algorithm that is a generalization of the project algorithm in the proof of \Cref{lem:embed-tree}. Roughly speaking, we keep decreasing each $d_i$ at the same rate, and whenever a triangle inequality is tight, we stop decreasing the corresponding $d_i$.

\begin{algorithm}[H]\label{alg:project}
\caption{Project a point to the metric}
\SetAlgoLined
$S \leftarrow [k]$.\;
\While{$S \neq \emptyset$}
{
    For any $i,j \in S$, $\Delta_{i,j} \leftarrow (d_i+d_j-\delta(t_i,t_j))/2$.\;
    For any $i\in S$ and $j \notin S$, $\Delta_{i,j} \leftarrow d_i+x_j-\delta(t_i,t_j)$.\;
    For any $i \in S$, $\Delta_i \leftarrow \min_j\{\Delta_{i,j}\}$.\;
    $\Delta \leftarrow \min_{i\in S} \Delta_i$.\;
    \For{$i$ such that $\Delta_i = \Delta$}{
        $x_i \leftarrow d_i - \Delta$.\;
        $S \leftarrow S - i$.\;
    }
}
Return $x=(x_1,x_2,\dots,x_k)$.\;
\end{algorithm}

We first prove the following property of the projection.    

\begin{claim} \label{clm:project}
For any index $i$, there is an index $j$ such that $x_i+x_j = \delta(t_i,t_j)$ and $d_i-x_i \ge d_j-x_j$.
\end{claim}

\begin{proof}
When an index $i$ is removed from $S$ and $x_i$ is fixed, we have $d_i-x_i = \Delta = \Delta_i \le \Delta_{i,j}$ for any $j$. Thus, for any index $j$ inside $S$ after this iteration, $\Delta_{i,j}$ will only increase. Thus for any $i$ and $j$, $\Delta_{i,j}$ does not decrease in each iteration, and so are $\Delta_i$ and $\Delta$.

For any index $i$, consider the iteration that $x_i$ is fixed, we have an index $j$ such that $\Delta_{i,j} = \Delta = d_i-x_i$. If $j \in S$, then $\Delta \le \Delta_j \le \Delta_{i,j} \le \Delta$, and thus $\Delta_j = \Delta$, which menas $x_j$ is also fixed in the same iteration, and so $d_i-x_i = \Delta_{i,j} = d_j-x_j$ and $x_i+x_j = d_i+d_j - 2\Delta_{i,j} = \delta(t_i,t_j)$. If $j \notin S$, then $x_i + x_j = d_i - \Delta_{i,j} + x_j = \delta(t_i,t_j)$. Furthermore, since $x_j$ is fixed before this iteration, we have $d_i-x_i = \Delta \ge d_j-x_j$ since $\Delta$ is not decreasing in each iteration.
\end{proof}

Now we prove that after projection, the distance between two point will not increase.

\begin{lemma} \label{lem:proj_dis}
    For any two vertex $u,v$ in the graph, suppose $x=f(u)$ and $y=f(v)$, $\norm{x-y}_{\infty} \le \ell(u,v)$.
\end{lemma}

\begin{proof}
    For any $i$, define $\Delta^u_i$ and $\Delta^v_i$ as $\ell(u,t_i) - x_i$ and $\ell(v,t_i) - y_i$. By definition of $\ell_{\infty}$ norm and symmetry, it is sufficient to prove that for any $i$, $x_i-y_i \le \ell(u,v)$. By definition, 
    $$x_i-y_i = (\ell(u,t_1) - \Delta^u_i) - (\ell(v,t_1) - \Delta^v_i) \le \ell(u,v) - \Delta^u_i + \Delta^v_i.$$ 
    By \Cref{clm:project}, there is an index $j$ such that $x_i+x_j = \ell(t_i,t_j)$ and $\Delta^u_i \ge \Delta^u_j$. We have 
    $$x_i-y_i \le (\ell(t_i,t_j) - x_j) - (\ell(t_i,t_j) - y_j) = y_j - x_j.$$
    and 
    $$
    y_j - x_j \le \ell(u,v) - \Delta^v_j + \Delta^u_j \le \ell(u,v) - \Delta^v_j + \Delta^u_i.
    $$
    If $\Delta^v_i \le \Delta^v_j$, we have $2(x_i-y_i) \le 2\ell(u,v) - \Delta^v_j + \Delta^v_i \le 2\ell(u,v)$. If $\Delta^v_i > \Delta^v_j$, the index $j$ satisfies $\Delta^u_i \ge \Delta^u_j$, $\Delta^v_i > \Delta^v_j$ and $x_i-y_i \le y_j-x_j$. 

    By the same analysis, regarding $y_j-x_j$, we can either prove $y_j-x_j \le \ell(u,v)$ or there exists an index $i'$ such that $\Delta^v_i > \Delta^v_j \ge \Delta^v_{i'}$, $\Delta^u_i \ge \Delta^u_j > \Delta^u_{i'}$ and $x_i-y_i \le y_j-x_j \le x_{i'}-y_{i'}$. And in the later case, we can continue using the same analysis on $x_{j'}-y_{j'}$. If we keep doing the same analysis, we either prove $x_i - y_i \le \ell(u,v)$ or construct an infinite list of index $i_1,i_2,\dots$ such that for any $p$, $\Delta^u_{i_p}>\Delta^u_{i_{p+1}}$ and $\Delta^v_{i_p}>\Delta^v_{i_{p+1}}$, which is not possible. 
\end{proof}

\subsection{Project the Points in Tight Span to the Steiner Nodes}
In this section, we define the steiner nodes. For any vector $x$ in the tight span $\TS(T)$, define $\Gamma(x)$ as the graph whose vertex set is $[k]$, and two vertex $i$ and $j$ is connected if $x_i+x_j = \ell(t_i,t_j)$. We say the vector $x$ is an extreme point if $\Gamma(x)$ is connected and not bipartite. It is clear that any two extreme point $x$ and $y$ has different $\Gamma(x)$ and $\Gamma(y)$. Therefore, there are finite number of extreme points. 

